\chapter{Character Model}

\begin{multicols}{2}

\section{Character Skills and Actions}

\textbf{Character Actions}: All character actions in the game involve skills in some way. They represent the character's ability to do things. The skill values are added to the dice roll to determine the chance of success for the actions.

\textbf{Action Points (AP):} Action Points are a resource used to represent time. Actions cost AP to perform in situations where time is an important resource. The amount of AP each character receives at the beginning of each round is 6.

\textbf{Skills and Proficiencies}: Skill values are determined by the sum of the skill bonus and the proficiency bonus. Each skill bonus only influences a single skill and can be trained individually. Proficiency bonuses, on the other hand, can affect multiple skills. The default values for all skills and proficiencies is 0.

There are 19 skills, as described in the table:

\end{multicols}

\begin{stattable}{Skills}{@{}llll@{}}
Skill & Proficiency & Talent & Trainable \\
\midrule
Health & Constitution & - & no\\
\midrule
\rowcolor{fallowtint}
Strike & Melee & Dexterity  & yes \\
\rowcolor{fallowtint}
Defend & Melee & Dexterity & yes \\
\rowcolor{fallowtint}
Grapple & Melee + STR-10 & - & no \\
\midrule
Accuracy & Ranged & Dexterity & yes\\
Reflex & Ranged + Awareness & Dexterity & no\\
\midrule
\rowcolor{fallowtint}
Balance & AGI-10 & Dexterity  & yes \\
\rowcolor{fallowtint}
Climb & AGI-10 & Dexterity  & yes \\
\rowcolor{fallowtint}
Swim & AGI-10 & Dexterity  & yes \\
\midrule
Prestidigitation & Dexterity & Dexterity & yes \\
Stealth & - & Dexterity & yes \\
\midrule
\rowcolor{fallowtint}
Cunning & Awareness & Intellect  & yes \\
\rowcolor{fallowtint}
Detection & Awareness & - & no \\
\rowcolor{fallowtint}
Exploration & Knowledge + Awareness & Intellect  & knowledge \\
\rowcolor{fallowtint}
Knowledge & 2x knowledge or 2x INT & Intellect  & knowledge \\
\midrule
% Combustion & Sorcery & Intellect & yes \\
% Electromagnetism & Sorcery & Intellect & yes \\
% Entropy & Sorcery & Intellect & yes \\
% Radiation & Sorcery & Intellect & yes \\
% Animancy & Sorcery & Intellect & yes \\
% Telepathy & Sorcery & Intellect & yes \\
% Biomancy & Sorcery & Intellect & yes \\
% Devotion & Devotion & Spirit & yes \\
Insight & Charisma & Intellect & yes \\
Deception & Charisma & Intellect & yes \\
Persuasion & Charisma & Intellect & yes \\
Will & SPI & - & no \\
\end{stattable}

\begin{multicols}{2}

Skills marked as trainable can be improved individually. The ones that are not can only get bonuses from their proficiency. 

Skills where trainable is "knowledge" mean that the skill value depends on knowledge levels.

\subsection{Proficiencies:}

The trainable proficiencies are as follows. Proficiencies and other elements marked with the magic tag belong to the magic system and are omitted in a no-magic game.

\textbf{Melee (DEX)}: General skill in melee combat. Uses DEX as talent. \\
\textbf{Ranged Combat (DEX)}: General skill in ranged combat. Uses DEX as talent. \\
\textbf{Awareness (DEX)}: The character's ability to pay attention to what matters and not get distracted. Uses DEX as talent.\\
\textbf{Charisma (INT)}: This is the baseline skill of communication and empathy, that is useful in any social situation. It includes both knowing what to communicate and knowing who to communicate it to in verbal and non verbal form. Uses INT as talent. \\
\textbf{\magic Sorcery (INT)}: General skill in casting spells. Uses INT as talent. \\
\textbf{\magic Devotion (SPI)}: Level of alignment with a deity. Uses SPI as talent.\\
\textbf{Conviction (SPI)}: Alignment of worldview and instinct. Uses SPI as talent. \\

\subsection{Levels}

The expected levels for skills and proficiencies range from 0 to 5, but can be increased even further.

\textbf{Level 0, Novice:} Knows only what is common sense\\
\textbf{Level 1, Apprentice:} Can do the basics\\
\textbf{Level 2, Amateur:} Masters the basics and has some advanced knowledge\\
\textbf{Level 3, Professional:} Works at a professional level\\
\textbf{Level 4, Expert:} Capable of reproducing almost anything in the field \\
\textbf{Level 5, Master:} Capable of producing results of the highest level.\\
\textbf{Level 6, Grandmaster:} Is pushing the boundaries of knowledge in their era.\\

\section{Talent and Learning}

\textbf{Gaining Experience (XP):} Improving skills and proficiencies and learning abilities requires training during free time, which yields experience based on training time. A regular training session for skills and proficiencies takes 10 days for 1 XP and costs a certain amount of money based on the difficulty level of what one wants to learn.

\textbf{Improving proficiencies}: The cost of XP to improve a proficiency is equal to $4+4*n$, where n is the proficiency level one wishes to achieve. If the talent is lower than the level one wants to reach, this value is multiplied by $2^x$, where x is the difference between n and the talent. 

\textbf{Improving skills:} The logic is exactly the same as proficiencies, and the cost to increase a skill is equal to $2+2*n$, which is half the cost to increase the proficiency.

A character has 4 talents that represent their potential for growth.

\textbf{Constitution (CON)}: Represents the character's general health and ability to adapt physically.

\textbf{Dexterity (DEX)}: Represents the character's motor skills, spacial intelligence and reaction time. Includes all skills related to athletics and combat. These are well explained in the chapter on combat.

\textbf{Intellect (INT)}: Represents memory, logical reasoning, and ability to learn complex things. Includes all knowledge, sorcery, most social skills, exploration skills and cunning in combat. 

\textbf{Spirit (SPI)}: Represents the character's resolve. It is the primary attribute for convictions and devotions, as well as for their ability to maintain morale.

\textbf{Improving Talents:} During character creation, the player must choose the character's talents in order of priority. After that, talents can only be increased by spending karma. The karma cost to increase a talent is equal to $4+4*n$, where n is the talent level one wishes to acquire. The maximum talent value is 5.

\textbf{Level Limits:} The values of talent and proficiency range from 0 to 5, and cannot be increased beyond that. Skills can be increased indefinitely, but they become exponentially more expensive since the maximum talent is 5.

\subsection{Karma}

\textbf{Karma:} is used to increase talents. It is received upon completing missions and their secondary objectives, as well as being rewarded for good gameplay.

The GM can reward players for loyalty to their character, creative use of the system, and participation in dramatic events.

\textbf{Completing missions:} difficult missions award more karma. The karma award should be given based on the difficulty of the mission relative to the character's abilities. It is possible to break down the mission into smaller objectives and award per objective attempted. Objective completion is not necessary, although it should generally award a bit more. It should be worth more to attempt a difficult task and fail than complete easy ones. Completing many easy ones may bring a lot of money, but harder ones yield more karma.

\textbf{Loyalty to the character:} This means the player acts in ways that the character would, rather than what the player thinks is best. This is a reward for going against practical advantages to favor character alignment.

\textbf{Creative use of the system:} This is an extra reward for finding new ways to be effective with your character using the system. Suggesting map features during a chase, coming up with good reasons why an NPC should help the team, suggesting items that may help solve some problem, interpreting a magical object use it creatively. The same trick should not be rewarded twice. This is a reward for creativity.

\textbf{Participation in dramatic events:} Speeches, discussions between characters, exploration, and conversations that show curiosity about the world, demonstration of the character's ethical and moral flaws, and comedic moments. This is a reward for playing a part in memorable moments.

The estimated Karma award is approximately two thirds for completing missions and one third for the rest. This ensures that good behavior gets rewarded and celebrated but characters stay roughly matched in progress and players don't start gaming the reward system. 

A recommended karma award budget per session is about 4 to 6 for a regular session where some objective gets cleared. This is about one talent upgrade every 3 to 5 sessions.

% Karma can also be subtracted as a penalty for acting against the principles of any of these rewards.

\section{Constitution}

Constitution governs the enhancement of attributes, health, and \magic mutation abilities.

\textbf{Health:} The character's ability to heal wounds, recover from exhaustion, and resist toxins. The use of health is described in the chapter on survival.

\textbf{Injury level (IL):} The measure of a character's physical integrity. More on this in the combat chapter.

\subsection{Attributes}

Attributes represent the physical characteristics of the character. They are divided into Strength (STR), Agility (AGI), and Stamina (STA). The common values for attributes for healthy young humans are 10.

\textbf{Attributes as skills:} sometimes, an attribute can be used as a skill. For example, strength for lifting weight, agility to outrunning someone and stamina for holding breath. In that case, use 2x(attribute -10) as the skill value.

\textbf{Improving attributes:} The XP cost to improve attributes is equal to the target attribute value divided by 2. However, if the average of the 3 attributes exceeds $10+CON-SM$, the cost is doubled. For every point higher, it increases by +100\% again. It is impossible to have any attribute higher than 20.

\textbf{Changing attributes:} It is possible to lose 1 point in one attribute to increase another, but this costs 2 XP.

\subsubsection{STR}

Represents the weight and muscular strength of a character. Characters with higher STR can hit harder, resist damage better, and perform more impressive feats of strength.

\textbf{ Natural Toughness (TGH)} define how much damage a character can withstand before receiving injuries. Their default value is 0.5 x STR xDM.  This can be interpreted as a larger body, tougher flesh or a stocky build.

\subsubsection{AGI}

Represents the character's athletic ability. Agility contributes to proficiency in swimming, climbing, and balance. Additionally, it defines movement speeds and gives bonuses to some combat moves.

\subsubsection{STA}

Represents the total breath and physical energy of the character. It is used to define tolerance to exhaustion. STA is used as a resource for all types of actions that require physical effort and also for bursts of action.

\subsection{Size}

Size is a characteristic that defines how much space a creature occupies on the grid and modifies the physical qualities of characters, weapons, and armor.

\textbf{Size interpretation:} Consider that the weight of a creature increases 10 times every 2 categories or 3 times for each category. A human is a creature of size 3 and weighs about 50kg. Therefore, a creature of size 1 would be a cat weighing 5kg, and a horse would be a creature of size 5, as it weighs about 500kg.

\textbf{Size and Space Occupation: } A creature of size 3 in human form occupies 1 square or hexagon. For every 2 size categories, a character on a square grid will increase their dimensions by 1 square, and a hexagonal grid will occupy 3 spaces (placing the token in the middle of 3), then 7. Other creatures may have different shapes.

Two creatures of size 3 can occupy the same space, but cannot end a turn like that. A creature can occupy the same space as another of two sizes higher than itself.

\textbf{Scaling a creature:} The strength of a creature increases with the cross-sectional area of a muscle, while weight increases with volume. Therefore, when increasing the size of a creature without changing its shape, it receives a penalty of -1 STR and -1 natural AGI.

\textbf{Standard Deflection (SD):} is the standard difficulty to hit a target, which is based solely on its size and movement. The larger the target, the easier it is to hit. The standard deflection is equal to -5-1xSM for stationary targets and -2-1xSM for moving targets. It is impossible for any type of penalty to reduce a character's deflection to a value lower than the SD.

\textbf{Damage Multiplier (DM):} multiplies all values derived from STR of a character and their appropriately scaled gear. It's TGH, armor's defenses, weapon damage, abilities, and many spells.

\textbf{Skill Modifier (SM):} is a modifier added to or subtracted from the character's values, in the proportions shown in the table.

\end{multicols}

\begin{stattable}{Skill Modifier (SM)}{@{}lc@{}}
Skill & Modifier \\
\midrule
Grapple & +5x \\
% Strength & +5x \\
\rowcolor{fallowtint}
Reflex & -1x \\
% Deflection & -1x \\
% Weapon Precision & -1x \\
SD & -1x \\
\rowcolor{fallowtint}
Stealth & -2x \\
Climbing & -2x \\
\end{stattable}

\begin{multicols}{2}

\textbf{Movement Multiplier (MM):} multiplies all movement speeds and the range of all weapons.

\textbf{Volume Multiplier (VM):} multiplies the weight of the character and items when scaling.

\bigskip

\end{multicols}

\begin{stattable}{Size}{@{}lccccccc@{}}
Size & 1 & 2 & 3 & 4 & 5 & 6 & 7 \\
\midrule
DM & 0.5 & 0.75 & 1 & 1.5 & 2 & 3 & 4 \\
SM & -2 & -1 & 0 & +1 & +2 & +3 & +4 \\
MM & 0.5 & 1 & 1 & 1.5 & 1.5 & 2 & 2.5 \\
VM & 0.1 & 0.33 & 1 & 3.33 & 10 & 33.3 & 100 \\
\end{stattable}

\begin{multicols}{2}


\section{Races and Age}

\subsection{Races}
Each playable race comes with its own bonuses to attributes and talents. The table below shows the playable races. In a no-magic game, it is up to the GM to decide which races are available.


\end{multicols}

\begin{stattable}{Races}{@{}lccccccc@{}}
Race & STR & AGI & STA & INT & SPI & DEX & CON \\
\midrule
Surbit & 1 & 2 & 2 & 0 & 0 & 1 & 1 \\
\rowcolor{fallowtint}
Dwarf & 1 & -1 & 1 & 0 & 1 & 0 & 0 \\
Human & 0 & 0 & 0 & 0 & 0 & 0 & 0 \\
\rowcolor{fallowtint}
Halkir & 0 & 2 & 1 & 0 & 0 & 1 & 0 \\
Erthal & 2 & 0 & 1 & 0 & 0 & 0 & 1 \\
\end{stattable}

\begin{multicols}{2}

\textbf{\large Surbit}
\begin{proplist}
\item Small-sized creature
\item Has the Fur ability
\item Has the Bear Nose ability
\item Has the Tail I ability
\item Has the Claws ability
\end{proplist}

\textbf{\large Dwarves}
\begin{proplist}
\item Has the Bat Ears ability
\item Has the Tireless ability
\end{proplist}

\textbf{\large Humans}
\begin{proplist}
\item Increases two talents by 1 at creation.
\item Can choose one from Tireless, Large Hands, and Beautiful.
\end{proplist}

\textbf{\large Erthal}
\begin{proplist}
\item Has the Large Hands ability
\item Has the Regenerative Reservoir ability
\end{proplist}

\textbf{\large Halkir}
\begin{proplist}
\item Has the Beautiful ability
\item Has the Infrared Vision ability
\end{proplist}


\subsection{Aging}

A character's age influences their attributes and ability to learn new things. A character can be created as a child, a teenager, an adult, or older. Each of these ages has its characteristics.

\begin{proplist}
    \item Child: -2 STR, talents start at 0 but cost half the karma. Considered a size 2 creature.
    \item Teenager: -2 STR, +1 to talents.
    \item Adult: Standard.
    \item Aging: Receives -1 to CON and also to all attributes at each stage.
\end{proplist}

The time required to advance in the age category can vary between races.

\end{multicols}

\begin{stattable}{Aging}{@{}lc@{}}
Stage & Human \\
\midrule
Child & 8-13 \\
\rowcolor{fallowtint}
Teenager & 14-18 \\
Adult & 19-39 \\
\rowcolor{fallowtint}
Mature & 40-59 \\
Old & 60-79 \\
\rowcolor{fallowtint}
Very old & 80-death \\
\end{stattable}

\begin{multicols}{2}


\section{Intellect}

Intellect measures the logical capacity, planning, and decision-making abilities of characters. The skills governed by intellect involve exploration, detection knowledge, sorcery, and cunning.

\textbf{Exploration:} It is used in survival situations to determine the ability to travel without getting lost and to find things in varied terrains. The exploration skill is based on the sum of detection proficiency and specific knowledge of the area.

\textbf{Sorcery:} Uses sorcery proficiency added to a related knowledge to cast spells. More details are in the combat section.

\textbf{Cunning:} It is the ability to make good decisions quickly and be aware of danger. More details are in the combat section.


\subsection{Knowledge}

Knowledges are used as requisite to learn abilities, to interact with features during exploration, to enable and improve several downtime activities, and as a general measure of the character's intelligence.

\textbf{Learning Knowledge}: The knowledge listed in the areas of knowledge is learned exactly the same way as any skill, using INT as the talent. The levels are also tracked in the same scale.

\textbf{Usage of knowledges:}
\begin{proplist}
    \item \textbf{Knowledge as Enabler:} Many actions require an amount of a certain knowledge to be able to be performed. Attempting to perform that without the necessary knowledge level leads to a penalty of -5 per missing knowledge level to whatever test was called.
    \item \textbf{Knowledge as a requirement:} Some abilities are learned automatically at certain levels of the areas of knowledge. Some abilities cannot be learned without the proper knowledge requirement.
    \item \textbf{Knowledge in magic:} There are three magic related areas of knowledge. They are used when casting and learning spells as a wizard. More details in the Sorcery chapter.
    \item \textbf{Knowledge as a skill:} Some actions use knowledge level as component for a skill level. 
\end{proplist}

\subsubsection{Areas of Knowledge}

The formal areas of knowledge in which the character can be trained are:

\begin{proplist}
    \item \magic \textbf{Alchemy:} Knowledge about the usage and creation of inorganic magical items. Necessary for alchemy spells.
    \item \magic \textbf{Animancy:} Knowledge about consciousness and the nature of sorcery. What it does in the body, how it manifests and affects the environment, how to manipulate it. Necessary for animancy spells.
    \item \magic \textbf{Biomancy:} Knowledge about biological structures. Necessary for biomancy magic.
    \item \magic \textbf{Shamanism:} Helps with finding spirits, diplomacy with them.
    
    \item \textbf{Building:} Knowledge about structures, their common forms, how to build or destroy them. Tunnels, bridges, buildings, siege engines. Helps with exploration inside buildings.
    \item \textbf{Smithing:} Knowledge about the manufacture and maintenance of weapons, armor and mechanisms. Includes tools, locks and siege weapons.
    \item \textbf{Chemistry:} Knowledge about chemicals and their craft. Includes bombs, acids, toxic substances.
 
    \item \textbf{Medicine:} can make autopsy, potions, poisons and everything from first aid to surgery.

    \item \textbf{Politics:} involves forgery, disguise, and acquiring local information about people.
    \item \textbf{Linguistics:} Knowledge about languages and communication. Helps wit deciphering, translating and learning new languages.
    
    \item \textbf{Geography:} Allows map-making. Helps traveling skill. Make tests to guess where resource events should be.
    \item \textbf{Survival:} knowledge about the wilderness. It includes actions like cooking, butchering animals, preparing campfires, making natural shelters, and crafting simple objects from materials found in nature.
    \item \textbf{Zoology:} Knowledge about animals. Helps with diplomacy with animals, hunting, finding weaknesses.
\end{proplist}
\noindent

% \item Botany: Knowledge about plants, their location, shape, magical properties, poisons, fruit-bearing capacity, ecology, etc.
% \item Physics: Knowledge about electricity, radiation, magnetism and their craft.

\subsection{Social}

Characters have skills that define how proficient they are in social situations. Charisma is the main proficiency that governs social skills.

There are four main skills relevant to social situations: Persuasion, Insight, Deception and Will.


\skill{ Persuasion }{ Charisma }{ This is the skill to be interesting and convincing to others. It includes the ability to know what to say to whom and when to get their attention and interest.

This skill includes eloquence, performance, flow of conversation, and the ability to make others feel have a good impression. It is used in social interactions to convince others of something, to get their attention, and to make them feel good about themselves. }

\skill{ Insight }{Charisma}{ This is the skill used to understand others' motivations and detect deception. It is used in social interactions to understand others' intentions and desires, as well as to detect when someone is lying, even when the logic makes sense. 

This can also be used to detect a person or creature's emotional state. It can also be used to understand what a person values and what they want out of an interaction. It detects if someone is doing something for a different reason than the obvious.

NPCs have tags on them that indicate what they value and what they want out of an interaction. Insight can be used to identify these tags to make it easier to manipulate them.}

\skill{ Deception }{Charisma}{ This is the skill used to lie and not be noticed. It is important for someone to make performances, impersonate someone, infiltrate environments without being noticed socially. This is the skill used to tell lies and not be noticed. }


\skill{ Will }{ SPI }{ Will is like the character's presence and its mental fortitude. A character with strong will is more resistant to emotional stress and to magical mind altering effects. 

In social situations, it is used to intimidate people and to resist being intimidated. In combat and exploration, it is used in morale tests to resist being afraid.

The will defense is normally used passively and without AP cost as a defense against mind-altering spells and effects.

\textbf{Active test:} Will can also be used actively by spending 3 AP to attempt to overcome fear, anger, confusion, and effects that attack will defense. They become free from the effects on a critical or success. The DL is the same as that which caused the effect.

The active will test can be made if the character understands that they are under magical influence or if they are acting against their own interests due to emotion. If the emotional influence was subtle enough, no active will test can be made.
}


\subsection{Contacts}

\textbf{Contact}: is a trusted character. Contacts can answer sensitive questions, offer better prices, perform favors, and can buy and sell illegal items. The type of contact will define what they do. The types of contacts are described in the chapter on social interactions.


\section{Spirit}

Spirit governs the character's resolve: their Will, morale, and convictions. It is also relevant for Sorcery.

\subsection{Limit Actions:}

Each character must choose a limit action.

\textbf{Coward:} is compelled to use their movement surge to run away and try to hide. If STA is insuficient, rest and then surge if possible.

\textbf{Violent:} is compelled to use their fight surge to inflict damage on someone. If STA is insuficient, rest and then surge if possible.

\textbf{Tanatosis:} Falls to the ground unconscious for 1 turn and wakes up the next.

\textbf{Abusive:} is compelled to spend 3 AP to insult allies, increasing the morale test DL by +1 for the team. If nobody is around, default to coward.

\textbf{Traitor:} is compelled to use their reaction surge to negotiate with the opponent. If the target is unable to communicate, default to coward.

\textbf{Martyr:} is compelled to use their reaction surge to protect an ally and reposition. If STA is insuficient, rest and then surge if possible.


\subsection{Convictions}

Convictions provide access to various buffs, exclusive abilities, and are the only way to increase the character's will, but they impose some restrictions on the character's behavior. 

Every odd level of conviction increases will by 1. A conviction at level 1 gives +1 will, a level 3 gives +2 and a level 5, +3.

\textbf{Acquiring a conviction:} A character must adopt one worldview conviction and one instinct conviction during character creation. They start at level 0. Each has a defined way of gaining XP through the character's actions.

\textbf{Evolving Convictions}: Each conviction has its XP and level tracked individually. SPI is the talent used to determine the experience cost of improving the conviction. Conviction abilities are automatically acquired at each level.

It is possible to lose conviction XP. If the session ends with negative XP, the proficiency loses a level. Each conviction has its condition for XP loss. If the level reaches 0, the conviction is considered renounced.

There are 8 convictions arranged along two axes: Extrospective - Introspective and Order - Chaos. 4 of them are worldview convictions and 4 are instinct convictions. The convictions are detailed in the chapter on skills.

\subsection{\magic Deities and Devotion}

The following is part of the magic system and is omitted in a no-magic game.

One can choose to be a cleric of a deity. To do this, one must have at least one conviction in common with the deity and acquire the cleric ability.

Being a cleric of a deity grants access to cleric abilities and also allows one to learn all the skills and spells favored by the deity using devotion XP and the devotion level as talent. The Devotion skill can be used to cast spells related to the entity.

Each deity has attitudes that are promoted and disapproved.

\subsubsection{Ateu}

Values hard work and commitment. \\

CONTRACT: Gains the Contract ability. When the contract is made, both characters define a completion condition and a penalty. Gains 1 XP for making a contract, a variable amount of XP for completing the contract, and loses the same amount of XP for failing to comply. The loss can be mitigated by paying a penalty. \\

CANCELLATION: Atheist clerics always know if the parties are alive and if they have upheld the contract, as long as they are within areas of divine reach. The contracted individuals are marked and can be tracked by any Atheist cleric until they pay the penalty or complete the contract. \\

Convictions: Ferocity and Adaptation \\

Abilities: Gas mask, magnetic shield, electric skin \\

Spells: Lamp, Revealing Light, Dazzling Flash, Sharpen, Rust, Repair metal, Start fire, Flame thrower, sustained lightning, activate explosive, alchemical grenade, ghostly sound, Taser, Firestorm, Aurora, freeze water, rain. \\

\subsubsection{Iskar}

Values courage and justice \\

TRIAL: Learns martial spirit and uses devotion instead of animancy as a skill. \\

REDEMPTION: Swears to redeem someone's soul. If able to subjugate them, allows the choice between the effects of mind disintegration as redemption or death. Gains variable XP each time this occurs or half that if the target does not accept and is killed. Loses half the value if entering combat and surrendering. \\

Convictions: Fatalism and Stoicism \\

Abilities: Synesthesia, Spiritual Strike, central intelligence, blind combat, psychoactive assistance, chaotic thorn, \\

\subsubsection{Minala}

Values 

FAITH TEST: Gains healing food. Heals if true and well-intentioned, but poisons and loses nutritional value if not. \\

RESURRECTION: Faithful clerics of Minala can be revived if placed inside one of her temples, following the restrictions of the resurrection ability. \\

Convictions: Divine Providence and Guardian \\

Abilities: \\

Spells: Accelerate Blooming, enhanced sleep, treehouse, control tree, healing touch, Defibrillate, regeneration, Reattach limb\\

\subsubsection{Rivka}

Values diplomacy and intelligence \\

CONFESSION: Gains the ability to alter memories. Receives 1 XP for using it. Loses 4 XP if revealing the secret or if the person realizes that their memory has been altered. Performs a confession ritual in which they see the memories associated with the story. Can alter the person's memory, but must not let them perceive the alteration. \\

DIVINE GUIDE: The cleric receives 1 exhaustion to create a ghost of Rivka. The ghost may help the group or not, depending on its will. \\

Convictions: Hedonism and Domination \\

Abilities: Synesthesia, read surface thoughts, central intelligence, sculpt face, universal communication, spiritual walk. \\

Spells: telepathic link, scare, provoke, calm, visual interference, suggestion, ghostly sound, Command, possession, create fog, rescue, psychoactive assistance, astral travel, torment, anchor soul, and Command. \\

% end --------------------------------------------------------------------------------------------------

\section{Character Creation}

\textbf{Step 1:} Choose the character focus. Wealth, DEX, CON, INT, or SPI. Assign each focus a priority level, from A to E. \\
It is possible to start at an older age and be allowed to select two focuses at A and none at E.

\textbf{Step 2: Race.} Choose a race, which grants racial abilities and bonuses to attributes and talents. Add the attribute bonuses after step 4. At this stage, you also choose the age. 

\textbf{Step 3: Wealth.} The character can buy things worth this amount, but not training. Only 10\% can be taken into gameplay. \\
A: 100 GP, \\
B: 60 GP, \\
C: 30 GP, \\
D: 10 GP, \\
E: 5 GP.\\

\textbf{Step 4: CON.} Characters start with 8 in all attributes and receive XP to increase them depending on the focus. \\
A. CON=2, +60 XP. \\
B. CON=1, +42 XP. \\
C. CON=1, +30 XP.\\
D. CON=0, +24 XP. \\
E. CON=0, +16 XP.\\

\textbf{Step 5: SPI.} Starts with XP to distribute among the proficiencies, skills, and SPI abilities. \\
A: SPI=2, 60 XP,  \\
B: SPI=1, 42 XP,  \\
C: SPI=1, 30 XP, \\
D: SPI=0, 24 XP, \\
E: SPI=0, 16 XP. \\

\textbf{Step 6: INT.} Starts with XP to distribute among the proficiencies, skills, and INT abilities. \\
A: INT=2, 60 XP \\
B: INT=1, 42 XP, \\
C: INT=1, 30 XP, \\
D: INT=0, 24 XP, \\
E: INT=0, 16 XP. \\
All characters start with one language at level 4.\\

\textbf{Step 7: DEX.} Starts with XP to distribute among the proficiencies, skills, and DEX abilities. \\
A: DEX=2, 60 XP. \\
B: DEX=1, 42 XP, \\
C: DEX=1, 30 XP \\
D: DEX=0, 24 XP \\
E: DEX=0, 16 XP.\\

\end{multicols}
